\subsection*{一步一步建立直观}

\paragraph{1. 先看小桨轮，而非单个粒子}

选择纯旋转场。每个点沿圆周运动，速度处处与半径垂直；把一个很小的桨轮放入场中，它会旋转，但一小团粒子的面积不会改变。对于场 $F=(-\omega y,\omega x)$，散度为零，平面旋度为 $2\omega$。旋度读数与角速度差一个因子二。



\paragraph{2. 再看小片区域是否膨胀}

选择径向增长场，速度与位置成正比。小片区域的两条边都按同一比例增长，面积的瞬时相对增长率因此是 $2\lambda$。这解释散度为什么把不同坐标方向的伸缩率相加。单条轨线远离原点，并不单独构成散度的定义。



\paragraph{3. 旋转与压缩可以同时发生}

第三种场把旋转与收缩叠加，粒子沿螺线向内运动。右图展示的是这一常系数场的轨线；散度与旋度分别检测局部面积变化和局部转动。一般空间变化的场需要在每个位置做局部计算，不能只看整条轨迹像不像旋涡。


